\documentclass[11pt,a4paper]{article}
\usepackage[T1]{fontenc}
\usepackage[utf8]{inputenc}
\usepackage[spanish]{babel}
\usepackage{amsmath,amssymb,mathtools,amsthm}
\usepackage{geometry,lmodern,microtype,array,booktabs,xcolor}
\usepackage[most]{tcolorbox}
\usepackage[hidelinks]{hyperref}
\geometry{margin=2.4cm}
\setlength{\parindent}{0pt}\setlength{\parskip}{.5em}
\newcommand{\N}{\mathbb{N}}\newcommand{\Z}{\mathbb{Z}}\newcommand{\Q}{\mathbb{Q}}\newcommand{\R}{\mathbb{R}}\newcommand{\C}{\mathbb{C}}
\newcommand{\ord}{\operatorname{ord}}\newcommand{\gen}[1]{\langle #1\rangle}
\newcommand{\Ker}{\operatorname{Ker}}\newcommand{\Imf}{\operatorname{Im}}\newcommand{\Aut}{\operatorname{Aut}}\newcommand{\mcd}{\operatorname{mcd}}\newcommand{\mcm}{\operatorname{mcm}}
\newtcolorbox{definicion}[1][]{colback=gray!6,colframe=gray!55,boxrule=.6pt,arc=2mm,left=2mm,right=2mm,top=1.5mm,bottom=1.5mm,breakable,#1}
\newtcolorbox{resultado}[1][]{colback=gray!4,colframe=black!45,boxrule=.6pt,arc=2mm,left=2mm,right=2mm,top=1.5mm,bottom=1.5mm,breakable,#1}
\newtcolorbox{ejemplo}[1][]{colback=gray!2,colframe=gray!45,boxrule=.5pt,arc=2mm,breakable,#1}
\newtcolorbox{demostracion}[1][]{colback=white,colframe=gray!35,boxrule=.5pt,arc=1mm,breakable,#1}
\begin{document}
\begin{center}{\LARGE\bfseries Estructuras algebraicas}\\[.4em]{\Large Ejemplos de grupos}\\[.4em]{\large Fichas de Propiedades y comprobaciones}\end{center}
\tableofcontents\newpage
\section{Grupos cíclicos y numéricos}

\subsection{$(\Z_4,+)$}
\begin{ejemplo}
Clases $\{0,1,2,3\}$, sumadas módulo $4$.
\end{ejemplo}
\textbf{Propiedades}
\begin{itemize}
\item Orden $4$.
\item Abeliano, cíclico, finito y soluble.
\item Neutro $0$; inverso de $x$ es $-x\pmod 4$.
\item Generadores: $1$ y $3$.
\item Todos sus subgrupos son normales.
\end{itemize}
\textbf{Comprobación o ejemplo.} $\ord(2)=2$, mientras que $\ord(1)=4$.

\subsection{$(\Z_n,+)$}
\begin{ejemplo}
Clases de congruencia módulo $n$, con suma módulo $n$.
\end{ejemplo}
\textbf{Propiedades}
\begin{itemize}
\item Orden $n$.
\item Abeliano, cíclico y soluble.
\item Neutro $\overline0$; inverso de $\overline a$ es $\overline{-a}$.
\item La clase $\overline k$ genera el grupo si y solo si $\mcd(k,n)=1$.
\item Todos sus subgrupos son normales.
\end{itemize}
\textbf{Comprobación o ejemplo.} $\Z_6$ tiene generadores $1$ y $5$, pero no $2$.

\subsection{$(\Z,+)$ y $n\Z$}
\begin{ejemplo}
Los enteros y sus subgrupos de múltiplos de $n$, con la suma.
\end{ejemplo}
\textbf{Propiedades}
\begin{itemize}
\item Abelianos y solubles.
\item Para $n\ne0$, $n\Z\cong\Z$: son cíclicos infinitos.
\item No son grupos de torsión ni divisibles.
\item Todos sus subgrupos son normales.
\end{itemize}
\textbf{Comprobación o ejemplo.} $6\Z\cap10\Z=30\Z$ y $6\Z+10\Z=2\Z$.

\subsection{$(\Q,+)$, $(\R,+)$ y $(\C,+)$}
\begin{ejemplo}
Los racionales, reales y complejos con suma.
\end{ejemplo}
\textbf{Propiedades}
\begin{itemize}
\item Abelianos, infinitos y solubles.
\item Son divisibles y libres de torsión.
\item El neutro es $0$; inverso de $x$ es $-x$.
\item No son cíclicos.
\end{itemize}
\textbf{Comprobación o ejemplo.} La ecuación $nx=q$ se resuelve tomando $x=q/n$.

\subsection{$(\Q^*,\cdot)$, $(\R^*,\cdot)$ y $(\C^*,\cdot)$}
\begin{ejemplo}
Los números distintos de cero con la multiplicación habitual.
\end{ejemplo}
\textbf{Propiedades}
\begin{itemize}
\item Abelianos, infinitos y solubles.
\item Neutro $1$; inverso de $x$ es $1/x$.
\item No son cíclicos.
\item Contienen torsión: $-1$ tiene orden $2$.
\end{itemize}
\textbf{Comprobación o ejemplo.} $0$ debe excluirse: no tiene inverso multiplicativo.

\subsection{$(\Z_n^*,\cdot)$}
\begin{ejemplo}
Unidades módulo $n$, con producto módulo $n$.
\end{ejemplo}
\textbf{Propiedades}
\begin{itemize}
\item Abeliano, finito y soluble.
\item Orden $\varphi(n)$ (función de Euler).
\item No siempre es cíclico.
\item Cada elemento tiene inverso multiplicativo módulo $n$.
\end{itemize}
\textbf{Comprobación o ejemplo.} $\Z_8^*=\{1,3,5,7\}\cong\Z_2\times\Z_2$, por lo que no es cíclico.

\subsection{$\Z_2\times\Z_2$}
\begin{ejemplo}
Pares de clases módulo $2$, sumados coordenada a coordenada.
\end{ejemplo}
\textbf{Propiedades}
\begin{itemize}
\item Orden $4$, abeliano y soluble.
\item No cíclico.
\item Todos los elementos no neutros tienen orden $2$.
\item Todos sus subgrupos son normales.
\end{itemize}
\textbf{Comprobación o ejemplo.} $(1,0)+(1,0)=(0,0)$ y ocurre lo mismo con $(0,1)$ y $(1,1)$.

\subsection{$\Z_m\times\Z_n$}
\begin{ejemplo}
Producto directo de dos grupos cíclicos, con suma coordenada a coordenada.
\end{ejemplo}
\textbf{Propiedades}
\begin{itemize}
\item Orden $mn$.
\item Abeliano y soluble.
\item Es cíclico precisamente si $\mcd(m,n)=1$.
\item El orden de $(a,b)$ es $\mcm(\ord(a),\ord(b))$.
\end{itemize}
\textbf{Comprobación o ejemplo.} $\Z_2\times\Z_3\cong\Z_6$; $\Z_2\times\Z_4$ no es cíclico.

\section{Simetrías y permutaciones}

\subsection{$\Delta_n$ (grupo diédrico)}
\begin{ejemplo}
Simetrías de un polígono regular de $n\ge3$ lados: $\Delta_n=\{r^k,sr^k:0\le k<n\}$ con $r^n=s^2=e$ y $srs=r^{-1}$.
\end{ejemplo}
\textbf{Propiedades}
\begin{itemize}
\item Orden $2n$.
\item No abeliano para $n\ge3$.
\item Generado por $r$ y $s$; no cíclico.
\item Todas las reflexiones tienen orden $2$.
\item El subgrupo de rotaciones $\gen r\cong\Z_n$ es normal.
\end{itemize}
\textbf{Comprobación o ejemplo.} En $\Delta_3$, $sr\ne rs$ y $\Delta_3\cong S_3$.

\subsection{$S_n$ (grupo simétrico)}
\begin{ejemplo}
Todas las permutaciones de $n$ símbolos con composición.
\end{ejemplo}
\textbf{Propiedades}
\begin{itemize}
\item Orden $n!$.
\item No abeliano si $n\ge3$; $S_1,S_2$ son abelianos.
\item Para $n\ge3$, no es cíclico.
\item Contiene el subgrupo normal $A_n$ de índice $2$.
\item Para $n\ge5$, no es soluble.
\end{itemize}
\textbf{Comprobación o ejemplo.} En $S_3$, $(1\,2)(2\,3)\ne(2\,3)(1\,2)$.

\subsection{$A_n$ (grupo alternado)}
\begin{ejemplo}
Permutaciones pares de $n$ símbolos.
\end{ejemplo}
\textbf{Propiedades}
\begin{itemize}
\item Orden $n!/2$ para $n\ge2$.
\item Normal en $S_n$.
\item Para $n\ge5$, simple, no abeliano y no soluble.
\item Para $n=3$, $A_3\cong\Z_3$ y es cíclico.
\end{itemize}
\textbf{Comprobación o ejemplo.} $A_4$ tiene $12$ elementos: identidad, $3$ dobles trasposiciones y $8$ ciclos de longitud $3$.

\subsection{$A_5$}
\begin{ejemplo}
Grupo alternado de cinco letras.
\end{ejemplo}
\textbf{Propiedades}
\begin{itemize}
\item Orden $60$.
\item No abeliano y no cíclico.
\item Simple y perfecto.
\item No soluble.
\end{itemize}
\textbf{Comprobación o ejemplo.} El único subgrupo normal propio de $A_5$ es el trivial.

\section{Grupos infinitos, matrices y cocientes}

\subsection{$\Q/\Z$}
\begin{ejemplo}
Clases de números racionales que difieren en un entero, con suma.
\end{ejemplo}
\textbf{Propiedades}
\begin{itemize}
\item Abeliano, infinito y soluble.
\item Grupo de torsión y divisible.
\item No es cíclico.
\end{itemize}
\textbf{Comprobación o ejemplo.} La clase $1/6+\Z$ tiene orden $6$.

\subsection{$\mu_\infty$}
\begin{ejemplo}
Todas las raíces complejas de la unidad, bajo multiplicación.
\end{ejemplo}
\textbf{Propiedades}
\begin{itemize}
\item Abeliano, infinito y de torsión.
\item Divisible; no cíclico.
\item Isomorfo a $\Q/\Z$.
\end{itemize}
\textbf{Comprobación o ejemplo.} $e^{2\pi i/5}$ tiene orden $5$.

\subsection{$GL_n(\R)$}
\begin{ejemplo}
Matrices reales invertibles $n\times n$, bajo multiplicación.
\end{ejemplo}
\textbf{Propiedades}
\begin{itemize}
\item Grupo infinito.
\item Para $n\ge2$, no abeliano ni cíclico.
\item Neutro $I_n$; inverso la matriz inversa.
\item La aplicación determinante toma valores en $\R^*$.
\end{itemize}
\textbf{Comprobación o ejemplo.} $AB$ y $BA$ pueden ser distintos incluso siendo $A,B$ invertibles.

\subsection{$SL_n(\R)$}
\begin{ejemplo}
Matrices reales de determinante $1$, con multiplicación.
\end{ejemplo}
\textbf{Propiedades}
\begin{itemize}
\item Subgrupo normal de $GL_n(\R)$ (núcleo del determinante).
\item Infinito para $n\ge2$.
\item No abeliano para $n\ge2$.
\end{itemize}
\textbf{Comprobación o ejemplo.} $\begin{pmatrix}1&1\\0&1\end{pmatrix}\in SL_2(\R)$.

\subsection{$F(S)$ (grupo libre)}
\begin{ejemplo}
Palabras reducidas formadas con generadores y sus inversos.
\end{ejemplo}
\textbf{Propiedades}
\begin{itemize}
\item Para dos o más generadores, no abeliano.
\item No tiene torsión no trivial.
\item El grupo libre con un generador es isomorfo a $\Z$.
\end{itemize}
\textbf{Comprobación o ejemplo.} En $F(a,b)$ las palabras $ab$ y $ba$ son distintas.

\subsection{$G/N$ (grupo cociente)}
\begin{ejemplo}
Clases laterales de un normal $N\trianglelefteq G$.
\end{ejemplo}
\textbf{Propiedades}
\begin{itemize}
\item Si $G$ es finito, $|G/N|=[G:N]=|G|/|N|$.
\item Neutro $N$; inverso de $gN$ es $g^{-1}N$.
\item Si $G$ es abeliano, el cociente también es abeliano.
\end{itemize}
\textbf{Comprobación o ejemplo.} $\Z_{12}/\gen4\cong\Z_4$ y $\Delta_3/\gen r\cong\Z_2$.

\section{Operaciones que no forman grupos}

\subsection{$(\N,+)$ y $(\Q,\cdot)$}
\begin{ejemplo}
Casos útiles para distinguir semigrupos, monoides y grupos.
\end{ejemplo}
\textbf{Propiedades}
\begin{itemize}
\item $(\N,+)$ es monoide si $0\in\N$, pero no grupo: faltan inversos aditivos.
\item $(\Q,\cdot)$ es monoide, pero no grupo: $0$ carece de inverso.
\item $(\Q^*,\cdot)$ sí es un grupo abeliano.
\end{itemize}
\textbf{Comprobación o ejemplo.} El fallo concreto en $\Q$ multiplicativo es $0\cdot x\ne1$ para todo $x$.

\subsection{Operaciones inventadas}
\begin{ejemplo}
En $\R\setminus\{-1\}$ se define $a*b=a+b+ab$; en $\R^*$ se define $a\star b=3ab$.
\end{ejemplo}
\textbf{Propiedades}
\begin{itemize}
\item Ambas operaciones forman grupos abelianos.
\item Para $*$: neutro $0$ e inverso $-a/(a+1)$.
\item Para $\star$: neutro $1/3$ e inverso $1/(9a)$.
\end{itemize}
\textbf{Comprobación o ejemplo.} $1+(a*b)=(1+a)(1+b)$ explica la primera operación mediante transporte de estructura.

\end{document}
